\documentclass[10pt]{beamer}
\usetheme{metropolis}

\usepackage{amsmath,amssymb}
\usepackage{mathtools}
\usepackage{xeCJK}
\setCJKmainfont{Adobe 楷体 Std}
\usepackage{tikz-cd}

% 1. 开启标题背景色填充
\metroset{block=fill}

% theorem environments
\newtheorem{proposition}{Proposition}
\newtheorem{exercise}{Exercise}

% notation
\newcommand{\R}{\mathbb{R}}
\newcommand{\Q}{\mathbb{Q}}
\newcommand{\C}{\mathbb{C}}
\newcommand{\F}{\mathcal{F}}
\newcommand{\RP}{\mathbb{R}\mathrm{P}}

\title{Ch1\quad Manifolds}
\subtitle{Differentiable Manifolds}
\date{September 17, 2026}
\author{LXJ}

\begin{document}

\maketitle

%====================================================================
\section{Notations}

\begin{frame}{Notations (I)}
	\begin{itemize}
		\item \textbf{injective}: A map $f: A \to B$; if $\forall\, x, y \in A$,
		      $x \neq y \Rightarrow f(x) \neq f(y)$, then $f$ is an
		      \emph{injective map}.
		\item \textbf{surjective}: A map $f: A \to B$ is a \emph{surjective map}
		      if $f$ is onto, i.e.\ $f(A) = B$.
		\item A \emph{commutative diagram} (of maps)
		      \[
			      \begin{tikzcd}[column sep=small, row sep=small, ampersand replacement=\&]
				      \& A \arrow[dl, "f"'] \arrow[dr, "h"] \& \\
				      B \arrow[rr, "g"'] \& \& C
			      \end{tikzcd}
			      \qquad\text{means}\qquad h = g \circ f .
		      \]
		\item $\R^d$ —— $d$-dim Euclidean space.
	\end{itemize}
\end{frame}

\begin{frame}{Notations (II)}
	\begin{itemize}
		\item The function $r_i: \R^d \to \R$, $(x_1, \dots, x_d) \mapsto x_i$,
		      is called the \emph{$i$-th (canonical) coordinate function} on
		      $\R^d$.\\
		      If $f: X \to \R^d$, then $f^i \coloneqq r_i \circ f$ is the
		      \emph{$i$-th component function} of $f$.
		\item $B_p(r)$ —— the \emph{open ball} of radius $r$ about $p$.
		\item Given a function $f: X \to \R$,
		      \[
			      \operatorname{supp} f \coloneqq \overline{f^{-1}(\R \setminus \{0\})}
		      \]
		      is called the \emph{support} of $f$.
		\item Given a multi-index $\alpha$:
		      \[
			      \frac{\partial^\alpha}{\partial x^\alpha}
			      \coloneqq
			      \frac{\partial^{|\alpha|}}{\partial x_1^{\alpha_1} \cdots \partial x_d^{\alpha_d}} .
		      \]
	\end{itemize}
\end{frame}

%====================================================================
\section{Differentiable Manifolds}

\begin{frame}{Functions of class $C^k$}
	\begin{definition}[$C^k$ functions]
		Let $U \subset \R^d$ be open, $f: U \to \R$. We say $f$ is
		\emph{differentiable of class $C^k$} on $U$ (for $k \geq 0$) if
		$\dfrac{\partial^\alpha}{\partial x^\alpha} f$ exists and is continuous
		for all $|\alpha| \leq k$.
		\begin{itemize}
			\item In particular, $f$ is $C^0$ if $f$ is continuous.
			\item $f: U \to \R^n$ is $C^k$ if each component of $f$ is $C^k$ on $U$.
			\item $f$ is $C^\infty$ if $f$ is $C^k$ for all $k \geq 0$.
		\end{itemize}
	\end{definition}
\end{frame}

\begin{frame}{Locally Euclidean spaces}
	\begin{definition}[Locally Euclidean space]
		A \emph{locally Euclidean space} $M$ of dim $d$ is a Hausdorff
		topological space $M$ for which each point $p \in M$ has a neighborhood
		homeomorphic to an open subset of $\R^d$.
		\begin{itemize}
			\item If $\varphi$ is a homeomorphism of a connected open set
			      $U \subset M$ to an open subset of $\R^d$, then $\varphi$ is
			      called a \emph{coordinate map};
			\item $x_i \coloneqq r_i \circ \varphi$ is called the
			      \emph{coordinate function};
			\item $(U, \varphi)$, i.e.\ $\bigl(U, (x_1, \dots, x_d)\bigr)$, is
			      called a \emph{coordinate system};
			\item if $p \in U$ s.t.\ $\varphi(p) = 0$, the coordinate system is
			      said to be \emph{centered at $p$}.
		\end{itemize}
	\end{definition}
\end{frame}

\begin{frame}{Differentiable structure}
	\begin{definition}[Differentiable structure]
		A \emph{differentiable structure} $\F$ of class $C^k$
		($1 \leq k \leq \infty$) on a locally Euclidean space $M$ is a
		collection of coordinate systems
		$\{(U_\alpha, \varphi_\alpha) \mid \alpha \in A\}$ satisfying:
		\begin{enumerate}
			\item[(a)] the coordinate systems \emph{cover} $M$:
			      $\displaystyle\bigcup_{\alpha \in A} U_\alpha = M$;
			\item[(b)] whenever $U_\alpha \cap U_\beta \neq \varnothing$, the map
			      $\varphi_\alpha \circ \varphi_\beta^{-1}$ defined on the image of
			      $U_\alpha \cap U_\beta$ is of class $C^k$,
			      $\forall\, \alpha, \beta \in A$;
			\item[(c)] the collection $\F$ is \emph{maximal}: if $(U, \varphi)$ is
			      a coordinate system s.t.\ $\varphi_\alpha \circ \varphi^{-1}$ and
			      $\varphi \circ \varphi_\alpha^{-1}$ are $C^k$ for all
			      $\alpha \in A$, then $(U, \varphi) \in \F$.
		\end{enumerate}
	\end{definition}
\end{frame}

\begin{frame}{Differentiable manifolds}
	\begin{definition}[Differentiable manifold]
		A \emph{$d$-dimensional differentiable manifold} of class $C^k$ is a
		pair $(M, \F)$ consisting of a second countable, locally Euclidean space
		$M$ of dim $d$, together with a differentiable structure $\F$ of class
		$C^k$.
	\end{definition}
	\begin{block}{Remark}
		If $\F_0 = \{(U_\alpha, \varphi_\alpha) \mid \alpha \in A\}$ is any
		collection of local coordinate systems satisfying (a) \& (b), there is a
		\emph{unique maximal} collection containing $\F_0$. Namely
		\[
			\F = \bigl\{ (U, \varphi) \bigm|
			\varphi \circ \varphi_\alpha^{-1} \ \&\ \varphi_\alpha \circ \varphi^{-1}
			\text{ are } C^k \text{ for all } (U_\alpha, \varphi_\alpha) \in \F_0 \bigr\}.
		\]
	\end{block}
	\begin{exercise}
		Prove $\F$ satisfies (a) \& (b), as well as (c).
	\end{exercise}
\end{frame}

\begin{frame}{Remarks}
	\begin{block}{Remark (other structures)}
		There are differentiable structures of class $C^\omega$ and
		\emph{complex analytic} structures: $C^\omega$ refers to the composition
		$\varphi_\alpha \circ \varphi_\beta^{-1}$ being given by convergent
		power series; a complex analytic structure refers to the composition
		given by holomorphic functions.
	\end{block}
	\begin{alertblock}{Remark}
		In this lecture, we only focus on the differentiable structure of class
		$C^\infty$.
	\end{alertblock}
\end{frame}

%====================================================================
\section{Examples of Differentiable Manifolds}

\begin{frame}{Examples (a), (b)}
	\begin{example}
		$\R^d$ is a manifold. A differentiable structure $\F$ is the
		\emph{maximal collection} of coordinate systems containing
		$(\R^d, \mathrm{Id})$.
	\end{example}
	\begin{example}
		A \emph{real} vector space $V$ has a natural manifold structure. Let
		$\{e_i\}$ be a basis of $V$ and $\{f_i\}$ the dual basis of $V^*$,
		i.e.\ $f_i(e_j) = \delta_{ij}$. Then $\bigl(V, (f_1, \dots, f_n)\bigr)$
		is a coordinate system.
	\end{example}
\end{frame}

\begin{frame}{Example (c): the $d$-sphere}
	\begin{example}[The $d$-sphere]
		$S^d = \{(x_1, \dots, x_{d+1}) \mid x_1^2 + \cdots + x_{d+1}^2 = 1\}$
		is a manifold. Denote
		\[
			n = (0, \dots, 0, 1), \qquad s = (0, \dots, 0, -1).
		\]
		Then $\bigl(S^d \setminus \{n\}, \pi_n\bigr)$,
		$\bigl(S^d \setminus \{s\}, \pi_s\bigr)$ are \emph{local coordinates}
		(stereographic projections) contained in its differentiable structure
		$\F$.
	\end{example}
	\begin{exercise}
		Verify that $\pi_n = \varphi_n^{-1}$, where
		\[
			\varphi_n(u) = \left( \frac{2u}{1 + \|u\|^2},\;
			\frac{\|u\|^2 - 1}{1 + \|u\|^2} \right) \in S^d,
			\qquad u = (u_1, \dots, u_d) \in \R^d .
		\]
	\end{exercise}
\end{frame}

\begin{frame}{Examples (d), (e)}
	\begin{example}
		An open subset $U$ of a differentiable manifold $(M, \F_M)$ is itself a
		differentiable manifold with
		\[
			\F_U = \bigl\{ \bigl(U_\alpha \cap U,\;
			\varphi_\alpha \big|_{U_\alpha \cap U}\bigr) \bigm|
			(U_\alpha, \varphi_\alpha) \in \F_M \bigr\}.
		\]
	\end{example}
	\begin{example}[The general linear group]
		$\mathrm{GL}_n(\R) \subset \R^{n^2}$ is an \emph{open subset} of
		$\R^{n^2}$, therefore it is a manifold --- since
		$\det: \R^{n^2} \to \R$ is continuous and
		\[
			\mathrm{GL}_n(\R) = \det{}^{-1}\bigl(\R \setminus \{0\}\bigr).
		\]
	\end{example}
\end{frame}

\begin{frame}{Example (f): product manifolds}
	\begin{example}[Product manifolds]
		Let $(M_1, \F_1)$ and $(M_2, \F_2)$ be differentiable manifolds of dim
		$d_1, d_2$. Then $M_1 \times M_2$ becomes a differentiable manifold of
		dim $d_1 + d_2$, with the differentiable structure $\F$ the
		\emph{maximal collection} containing
		\[
			\bigl\{ \bigl(U_\alpha \times V_\beta,\; (\varphi_\alpha, \psi_\beta)\bigr)
			\bigm| (U_\alpha, \varphi_\alpha) \in \F_1,\;
			(V_\beta, \psi_\beta) \in \F_2 \bigr\}.
		\]
	\end{example}
\end{frame}

\begin{frame}{Example (g): projective space $\RP^n$}
	\begin{example}[Projective space]
		Define $\RP^n$ to be the \emph{quotient topological space} of
		$\R^{n+1} \setminus \{0\}$:
		\[
			\RP^n = \bigl(\R^{n+1} \setminus \{0\}\bigr) \big/ {\sim},
		\]
		with ``$\sim$'' the equivalence relation
		\[
			X \sim Y \iff \exists\, \lambda \neq 0 \ \text{s.t.}\ X = \lambda Y .
		\]
		Introducing \emph{homogeneous coordinates} to denote points of $\RP^n$:
		\[
			[x] \coloneqq (x_1 : x_2 : \cdots : x_{n+1}).
		\]
	\end{example}
\end{frame}

\begin{frame}{Example (g): local coordinates on $\RP^n$}
	\begin{example}[Projective space (continued)]
		Consider the open set
		\[
			U_k = \bigl\{ x \in \R^{n+1} \mid x_k \neq 0 \bigr\} \big/ {\sim}
			\ \subset \RP^n ,
		\]
		and denote
		\[
			\varphi_k([x]) = \left( \frac{x_1}{x_k}, \dots,
			\widehat{\frac{x_k}{x_k}}, \dots, \frac{x_{n+1}}{x_k} \right) \in \R^n .
		\]
		Then $(U_k, \varphi_k)$ gives a local coordinate system of $\RP^n$.\\
		$\Rightarrow$ $\RP^n$ is a differentiable manifold with differentiable
		structure $\F$ containing
		\[
			\bigl\{ (U_k, \varphi_k) \bigm| k = 1, \dots, n+1 \bigr\}.
		\]
	\end{example}
\end{frame}

%====================================================================
\section{Smooth Functions and Maps}

\begin{frame}{$C^\infty$ functions and $C^\infty$ maps}
	\begin{definition}[$C^\infty$ functions and $C^\infty$ maps]
		We say $f: M \to \R$ is a \emph{$C^\infty$ function} on $M$ if
		$f \circ \varphi_\alpha^{-1}$ is $C^\infty$ for each coordinate map
		$(U_\alpha, \varphi_\alpha)$ on $M$.
		\par\medskip
		A continuous map $f: M \to N$ is a \emph{$C^\infty$ map} between
		manifolds $M$ and $N$ if
		\[
			\psi_\beta \circ f \circ \varphi_\alpha^{-1} \in C^\infty
			\qquad \text{for all } (U_\alpha, \varphi_\alpha)
			\text{ on } M \text{ and } (V_\beta, \psi_\beta) \text{ on } N .
		\]
	\end{definition}
\end{frame}

\end{document}
